% ---- score2 standard preamble (from: python3 score2_tools.py preamble) ----
\documentclass[10pt]{extarticle}
\usepackage[a4paper,margin=1.5cm]{geometry}
\usepackage[T1]{fontenc}
\usepackage{amsmath,amssymb}
\usepackage{graphicx}
\usepackage[export]{adjustbox}
\usepackage{multicol}
\usepackage{enumitem}
\usepackage{booktabs,array,tabularx}
\setlength{\columnsep}{0.7cm}
\setlength{\parindent}{0pt}
\setlength{\parskip}{3pt}
\pagestyle{empty}
\makeatletter
% figure registry (figdefs.tex is generated by the extract step)
\newcommand{\qfigdef}[4]{\@namedef{qf@file@#1}{#2}\@namedef{qf@w@#1}{#3}\@namedef{qf@wide@#1}{#4}}
\newcommand{\qf@check}[1]{\@ifundefined{qf@file@#1}{\PackageError{score2}{Unknown figure id #1}{}}{}}
\newcommand{\qf@img}[4][]{\includegraphics[width=#2\dimexpr\@nameuse{qf@w@#4}pt\relax,max width=#3,max height=0.45\textheight,#1]{\@nameuse{qf@file@#4}}}
% Full-width questions: finalize wraps every question that contains a wide figure in \qfullbegin...\qfullend,
% so the WHOLE question (stem, figures, options) is set across the page in reading order, then the two columns
% resume. The question's head - its number, stem and first figure (\qfullsplit marks the end of the head) -
% is kept on one page, so a number is never left at a page foot with its figure overleaf; the rest may break
% across pages normally, so no big gaps are left. Without a \qfullsplit, the whole question is kept together
% if it is under 45% of a page. Never write these yourself.
\newif\ifqf@full
\newif\ifqf@boxed
\newsavebox{\qf@box}
\newcommand{\qf@open}{\global\qf@fulltrue\par\global\qf@boxedtrue
  \setbox\qf@box\vbox\bgroup\hsize\textwidth\linewidth\textwidth}
\newcommand{\qf@emit}[1]{\par\egroup\global\qf@boxedfalse
  \ifdim\dimexpr\ht\qf@box+\dp\qf@box\relax<#1\textheight\noindent\usebox{\qf@box}\par\else\unvbox\qf@box\fi}
\newcommand{\qfullsplit}{\ifqf@boxed\qf@emit{0.85}\fi}
\newcommand{\qf@close}{\ifqf@boxed\qf@emit{0.45}\fi\global\qf@fullfalse}
\newcommand{\qfullbegin}{\end{multicols}\qf@open}
\newcommand{\qfullend}{\qf@close\begin{multicols}{2}}
\newcommand{\qfullstart}{\qf@open}
\newcommand{\qfulldone}{\qf@close}
% \qfig[scale]{id}: block figure, centred, directly where it is written (size set by the tool).
% Fallback if a wide figure is not inside a full-width question: break out of the columns in place - no floats.
\newcommand{\qfig}[2][1]{\qf@check{#2}%
  \ifqf@full\qf@blk{#1}{#2}\else
  \ifnum\@nameuse{qf@wide@#2}=1 \expandafter\@firstoftwo\else\expandafter\@secondoftwo\fi
  {\end{multicols}\begin{center}\qf@img{#1}{\textwidth}{#2}\end{center}\begin{multicols}{2}}%
  {\qf@blk{#1}{#2}}\fi}
\newcommand{\qf@blk}[2]{\par\smallskip{\centering\qf@img{#1}{\linewidth}{#2}\par}\smallskip}
% \qfigi[scale]{id}: inline figure (inside option lists / tables / minipages); never breaks columns.
\newcommand{\qfigi}[2][1]{\qf@check{#2}\qf@img[valign=c]{#1}{\linewidth}{#2}}
\makeatother
% structure macros - lint reads these to check completeness and figure placement
\newcommand{\papersec}[1]{\par\bigskip{\centering\Large\bfseries PAPER-#1\par}\medskip}
\newcommand{\partsec}[2]{\par\medskip\noindent{\large\bfseries PART-#1: #2}\par\smallskip}
\newcommand{\sect}[1]{\par\smallskip\noindent\textbf{#1}\par}
\newcommand{\qn}[1]{\par\smallskip\noindent\textbf{#1.}~}
\newcommand{\paraq}[2]{\par\smallskip\noindent\textbf{Paragraph for Questions #1 and #2}\par}
\newenvironment{qoptions}{\begin{itemize}[leftmargin=1.8em,itemsep=1pt,topsep=2pt,label={}]}{\end{itemize}}
\newcommand{\opt}[2]{\item[(#1)] #2}
% figure registry (generated by score2_tools.py finalize)
\qfigdef{p1_phy_s1_q01_1}{JEE Advanced (2021 Series) - Test 14 (06 Jun 2021)-img/p1_phy_s1_q01_1.png}{121.0}{0}
\qfigdef{p1_phy_s1_q02_1}{JEE Advanced (2021 Series) - Test 14 (06 Jun 2021)-img/p1_phy_s1_q02_1.png}{95.1}{0}
\qfigdef{p1_phy_s1_q06_1}{JEE Advanced (2021 Series) - Test 14 (06 Jun 2021)-img/p1_phy_s1_q06_1.png}{105.1}{0}
\qfigdef{p1_phy_s1_q06_A}{JEE Advanced (2021 Series) - Test 14 (06 Jun 2021)-img/p1_phy_s1_q06_A.png}{94.3}{0}
\qfigdef{p1_phy_s1_q06_B}{JEE Advanced (2021 Series) - Test 14 (06 Jun 2021)-img/p1_phy_s1_q06_B.png}{95.0}{0}
\qfigdef{p1_phy_s1_q06_C}{JEE Advanced (2021 Series) - Test 14 (06 Jun 2021)-img/p1_phy_s1_q06_C.png}{93.7}{0}
\qfigdef{p1_phy_s1_q06_D}{JEE Advanced (2021 Series) - Test 14 (06 Jun 2021)-img/p1_phy_s1_q06_D.png}{93.9}{0}
\qfigdef{p1_phy_s1_q07_1}{JEE Advanced (2021 Series) - Test 14 (06 Jun 2021)-img/p1_phy_s1_q07_1.png}{80.9}{0}
\qfigdef{p1_phy_s1_q08_1}{JEE Advanced (2021 Series) - Test 14 (06 Jun 2021)-img/p1_phy_s1_q08_1.png}{66.3}{0}
\qfigdef{p1_phy_s1_q09_1}{JEE Advanced (2021 Series) - Test 14 (06 Jun 2021)-img/p1_phy_s1_q09_1.png}{96.3}{0}
\qfigdef{p1_phy_s1_q10_1}{JEE Advanced (2021 Series) - Test 14 (06 Jun 2021)-img/p1_phy_s1_q10_1.png}{80.6}{0}
\qfigdef{p1_phy_s2_q01_1}{JEE Advanced (2021 Series) - Test 14 (06 Jun 2021)-img/p1_phy_s2_q01_1.png}{83.8}{0}
\qfigdef{p1_phy_s2_q02_1}{JEE Advanced (2021 Series) - Test 14 (06 Jun 2021)-img/p1_phy_s2_q02_1.png}{61.5}{0}
\qfigdef{p1_phy_s2_q04_1}{JEE Advanced (2021 Series) - Test 14 (06 Jun 2021)-img/p1_phy_s2_q04_1.png}{65.2}{0}
\qfigdef{p1_phy_s2_q05_1}{JEE Advanced (2021 Series) - Test 14 (06 Jun 2021)-img/p1_phy_s2_q05_1.png}{89.6}{0}
\qfigdef{p1_phy_s3_q01_1}{JEE Advanced (2021 Series) - Test 14 (06 Jun 2021)-img/p1_phy_s3_q01_1.png}{117.5}{0}
\qfigdef{p1_phy_s3_q02_1}{JEE Advanced (2021 Series) - Test 14 (06 Jun 2021)-img/p1_phy_s3_q02_1.png}{55.3}{0}
\qfigdef{p1_phy_s3_q04_1}{JEE Advanced (2021 Series) - Test 14 (06 Jun 2021)-img/p1_phy_s3_q04_1.png}{125.2}{0}
\qfigdef{p1_chem_s1_q05_1}{JEE Advanced (2021 Series) - Test 14 (06 Jun 2021)-img/p1_chem_s1_q05_1.png}{241.2}{0}
\qfigdef{p1_chem_s1_q05_A}{JEE Advanced (2021 Series) - Test 14 (06 Jun 2021)-img/p1_chem_s1_q05_A.png}{39.0}{0}
\qfigdef{p1_chem_s1_q05_B}{JEE Advanced (2021 Series) - Test 14 (06 Jun 2021)-img/p1_chem_s1_q05_B.png}{58.8}{0}
\qfigdef{p1_chem_s1_q05_C}{JEE Advanced (2021 Series) - Test 14 (06 Jun 2021)-img/p1_chem_s1_q05_C.png}{40.4}{0}
\qfigdef{p1_chem_s1_q05_D}{JEE Advanced (2021 Series) - Test 14 (06 Jun 2021)-img/p1_chem_s1_q05_D.png}{64.7}{0}
\qfigdef{p1_chem_s1_q06_1}{JEE Advanced (2021 Series) - Test 14 (06 Jun 2021)-img/p1_chem_s1_q06_1.png}{85.0}{0}
\qfigdef{p1_chem_s1_q06_A}{JEE Advanced (2021 Series) - Test 14 (06 Jun 2021)-img/p1_chem_s1_q06_A.png}{38.7}{0}
\qfigdef{p1_chem_s1_q06_B}{JEE Advanced (2021 Series) - Test 14 (06 Jun 2021)-img/p1_chem_s1_q06_B.png}{40.4}{0}
\qfigdef{p1_chem_s1_q06_C}{JEE Advanced (2021 Series) - Test 14 (06 Jun 2021)-img/p1_chem_s1_q06_C.png}{22.8}{0}
\qfigdef{p1_chem_s1_q06_D}{JEE Advanced (2021 Series) - Test 14 (06 Jun 2021)-img/p1_chem_s1_q06_D.png}{25.0}{0}
\qfigdef{p1_chem_s1_q07_1}{JEE Advanced (2021 Series) - Test 14 (06 Jun 2021)-img/p1_chem_s1_q07_1.png}{106.3}{0}
\qfigdef{p1_chem_s1_q09_A}{JEE Advanced (2021 Series) - Test 14 (06 Jun 2021)-img/p1_chem_s1_q09_A.png}{50.0}{0}
\qfigdef{p1_chem_s1_q09_B}{JEE Advanced (2021 Series) - Test 14 (06 Jun 2021)-img/p1_chem_s1_q09_B.png}{49.2}{0}
\qfigdef{p1_chem_s1_q09_C}{JEE Advanced (2021 Series) - Test 14 (06 Jun 2021)-img/p1_chem_s1_q09_C.png}{65.1}{0}
\qfigdef{p1_chem_s1_q09_D}{JEE Advanced (2021 Series) - Test 14 (06 Jun 2021)-img/p1_chem_s1_q09_D.png}{69.2}{0}
\qfigdef{p1_chem_s2_q03_A}{JEE Advanced (2021 Series) - Test 14 (06 Jun 2021)-img/p1_chem_s2_q03_A.png}{128.5}{0}
\qfigdef{p1_chem_s2_q03_B}{JEE Advanced (2021 Series) - Test 14 (06 Jun 2021)-img/p1_chem_s2_q03_B.png}{119.0}{0}
\qfigdef{p1_chem_s3_q02_1}{JEE Advanced (2021 Series) - Test 14 (06 Jun 2021)-img/p1_chem_s3_q02_1.png}{126.3}{0}
\qfigdef{p1_chem_s3_q03_1}{JEE Advanced (2021 Series) - Test 14 (06 Jun 2021)-img/p1_chem_s3_q03_1.png}{119.6}{0}
\qfigdef{p2_phy_s1_q01_1}{JEE Advanced (2021 Series) - Test 14 (06 Jun 2021)-img/p2_phy_s1_q01_1.png}{87.2}{0}
\qfigdef{p2_phy_s1_q02_1}{JEE Advanced (2021 Series) - Test 14 (06 Jun 2021)-img/p2_phy_s1_q02_1.png}{99.5}{0}
\qfigdef{p2_phy_s1_q04_1}{JEE Advanced (2021 Series) - Test 14 (06 Jun 2021)-img/p2_phy_s1_q04_1.png}{114.8}{0}
\qfigdef{p2_phy_s2_para1_2_1}{JEE Advanced (2021 Series) - Test 14 (06 Jun 2021)-img/p2_phy_s2_para1_2_1.png}{92.3}{0}
\qfigdef{p2_phy_s2_para3_4_1}{JEE Advanced (2021 Series) - Test 14 (06 Jun 2021)-img/p2_phy_s2_para3_4_1.png}{105.1}{0}
\qfigdef{p2_phy_s2_para5_6_1}{JEE Advanced (2021 Series) - Test 14 (06 Jun 2021)-img/p2_phy_s2_para5_6_1.png}{177.3}{0}
\qfigdef{p2_phy_s3_q02_1}{JEE Advanced (2021 Series) - Test 14 (06 Jun 2021)-img/p2_phy_s3_q02_1.png}{91.8}{0}
\qfigdef{p2_phy_s3_q03_1}{JEE Advanced (2021 Series) - Test 14 (06 Jun 2021)-img/p2_phy_s3_q03_1.png}{130.1}{0}
\qfigdef{p2_phy_s3_q04_1}{JEE Advanced (2021 Series) - Test 14 (06 Jun 2021)-img/p2_phy_s3_q04_1.png}{183.5}{0}
\qfigdef{p2_phy_s3_q05_1}{JEE Advanced (2021 Series) - Test 14 (06 Jun 2021)-img/p2_phy_s3_q05_1.png}{112.3}{0}
\qfigdef{p2_phy_s3_q06_1}{JEE Advanced (2021 Series) - Test 14 (06 Jun 2021)-img/p2_phy_s3_q06_1.png}{128.0}{0}
\qfigdef{p2_phy_s3_q07_1}{JEE Advanced (2021 Series) - Test 14 (06 Jun 2021)-img/p2_phy_s3_q07_1.png}{95.8}{0}
\qfigdef{p2_phy_s3_q08_1}{JEE Advanced (2021 Series) - Test 14 (06 Jun 2021)-img/p2_phy_s3_q08_1.png}{130.1}{0}
\qfigdef{p2_phy_s3_q10_1}{JEE Advanced (2021 Series) - Test 14 (06 Jun 2021)-img/p2_phy_s3_q10_1.png}{126.4}{0}
\qfigdef{p2_chem_s1_q03_D}{JEE Advanced (2021 Series) - Test 14 (06 Jun 2021)-img/p2_chem_s1_q03_D.png}{130.0}{0}
\qfigdef{p2_chem_s1_q04_A}{JEE Advanced (2021 Series) - Test 14 (06 Jun 2021)-img/p2_chem_s1_q04_A.png}{155.6}{0}
\qfigdef{p2_chem_s1_q04_2}{JEE Advanced (2021 Series) - Test 14 (06 Jun 2021)-img/p2_chem_s1_q04_2.png}{46.0}{0}
\qfigdef{p2_chem_s2_q05_A}{JEE Advanced (2021 Series) - Test 14 (06 Jun 2021)-img/p2_chem_s2_q05_A.png}{56.5}{0}
\qfigdef{p2_chem_s2_q05_B}{JEE Advanced (2021 Series) - Test 14 (06 Jun 2021)-img/p2_chem_s2_q05_B.png}{55.8}{0}
\qfigdef{p2_chem_s2_q05_C}{JEE Advanced (2021 Series) - Test 14 (06 Jun 2021)-img/p2_chem_s2_q05_C.png}{55.4}{0}
\qfigdef{p2_chem_s2_q05_D}{JEE Advanced (2021 Series) - Test 14 (06 Jun 2021)-img/p2_chem_s2_q05_D.png}{51.2}{0}
\qfigdef{p2_chem_s2_q06_A}{JEE Advanced (2021 Series) - Test 14 (06 Jun 2021)-img/p2_chem_s2_q06_A.png}{205.7}{0}
\qfigdef{p2_chem_s2_q06_B}{JEE Advanced (2021 Series) - Test 14 (06 Jun 2021)-img/p2_chem_s2_q06_B.png}{223.0}{0}
\qfigdef{p2_chem_s2_q06_C}{JEE Advanced (2021 Series) - Test 14 (06 Jun 2021)-img/p2_chem_s2_q06_C.png}{197.3}{0}
\qfigdef{p2_chem_s2_q06_D}{JEE Advanced (2021 Series) - Test 14 (06 Jun 2021)-img/p2_chem_s2_q06_D.png}{224.9}{0}
\qfigdef{p2_chem_s3_q03_1}{JEE Advanced (2021 Series) - Test 14 (06 Jun 2021)-img/p2_chem_s3_q03_1.png}{100.9}{0}
\qfigdef{p2_chem_s3_q09_1}{JEE Advanced (2021 Series) - Test 14 (06 Jun 2021)-img/p2_chem_s3_q09_1.png}{242.0}{0}
\qfigdef{p2_chem_s3_q10_1}{JEE Advanced (2021 Series) - Test 14 (06 Jun 2021)-img/p2_chem_s3_q10_1.png}{125.6}{0}
% ---- body skeleton ----
% \begin{document}
% \begin{center}{\Large\bfseries JEE (Main+Advanced) Score-II --- Test NN (DD Mon YYYY)}\end{center}
% \begin{multicols}{2}
%   \papersec{1}   (only for two-paper tests)
%   \partsec{1}{PHYSICS} \sect{Section-I (i): Single Correct} \qn{1} ... \qfig{p1_phy_q01_1} ...
% \end{multicols}
% \end{document}

\begin{document}
\begin{center}{\Large\bfseries JEE Advanced (2021 Series) --- Test 14 (06 Jun 2021)}\end{center}
\begin{multicols}{2}
\papersec{1}
\partsec{1}{PHYSICS}
\sect{SEC--1}
\qn{1} Precession is the result of the angular velocity of rotation (spin) and the angular velocity produced by the torque. It is an angular velocity about a line that makes an angle $\theta$ with the permanent rotation axis. In the diagram, a top of mass $m$ is performing precession motion. The torque due to gravitational force $F_g$ causes a change in the angular momentum $L$ in the direction of that torque causing the top to process with angular velocity $\omega_P$ which is given by: (Given $r$ is the distance of COM from point of rotation `P'. $I_s$ \& $\omega_s$ are the moment of inertia and angular velocity of top about spin axis respectively.)
\qfig{p1_phy_s1_q01_1}
\begin{qoptions}
\opt{A}{$\dfrac{mgr}{I_s\omega_s^2}$}
\opt{B}{$\dfrac{mgr^2}{I_s\omega_s^2}$}
\opt{C}{$\dfrac{\omega_sI_s}{mgr}$}
\opt{D}{$\dfrac{mgr}{\omega_sI_s}$}
\end{qoptions}
\qn{2} A spherical metal ball of radius `$r$' is lying at the bottom of stationary container containing liquid of density $\rho$ as shown in the figure. The force exerted on the upper hemispherical portion of the sphere due to pressure ($p_0$ = atmospheric pressure) is:
\qfig{p1_phy_s1_q02_1}
\begin{qoptions}
\opt{A}{$\dfrac{\pi r^2}3[3P_0+7r\rho g]$}
\opt{B}{$\dfrac{\pi r^2}2[3P_0+7r\rho g]$}
\opt{C}{$\pi r^2[3P_0+7r\rho g]$}
\opt{D}{$2\pi r^2[3P_0+7r\rho g]$}
\end{qoptions}
\qn{3} A real object is placed 1 cm above the optical axis of a convex lens of focal length 40cm. The object distance is 60cm. If the object now starts moving perpendicularly away from the optical axis with a speed = 10 cm/s, the speed of the image is
\begin{qoptions}
\opt{A}{5 cm/sec}
\opt{B}{10 cm/sec}
\opt{C}{20 cm/sec}
\opt{D}{40 cm/sec}
\end{qoptions}
\qn{4} In the Bohr's atomic model, electron revolves in an orbit with speed $v$. It produces magnetic field $B$ at centre. The magnetic dipole moment of the electron is $M$. Which of the following quantity is independent of quantum number of the orbit?
\begin{qoptions}
\opt{A}{$\dfrac{Mr}{Bv}$}
\opt{B}{$\dfrac{Br^2}{Mv^2}$}
\opt{C}{$\dfrac{Mr^2}{Bv^2}$}
\opt{D}{$\dfrac{Mv^2}{Br^3}$}
\end{qoptions}
\qn{5} Two identical blocks of same metal are at $20^\circ$C and $80^\circ$C respectively. Specific heat of the material of the two blocks increase with temperature. Which of the following is true about the final temperature $T_f$ when the two blocks are brought into contact? (Assume that no heat is lost to the surroundings):
\begin{qoptions}
\opt{A}{$T_f$ will be $50^\circ$C}
\opt{B}{$T_f$ will be more than $50^\circ$C}
\opt{C}{$T_f$ will be less than $50^\circ$C}
\opt{D}{$T_f$ can be either more than or less than $50^\circ$C depending on the precise variation of the specific heat with temperature.}
\end{qoptions}
\qn{6} A broad metal plate is connected to earth through a conducting wire as shown. An electron flies with constant velocity along a straight line above the plate at a distance much less than the linear dimensions of the plate. If the current $I$ flowing from the earth to the plate is considered to be positive, then which of the following graph correctly depicts the variation of current $I$ with time?
\qfig{p1_phy_s1_q06_1}
\begin{qoptions}
\opt{A}{\qfigi{p1_phy_s1_q06_A}}
\opt{B}{\qfigi{p1_phy_s1_q06_B}}
\opt{C}{\qfigi{p1_phy_s1_q06_C}}
\opt{D}{\qfigi{p1_phy_s1_q06_D}}
\end{qoptions}
\qn{7} An electric circuit consists of a battery emf $E=110$V, with internal resistance $5\Omega$ and two resistors connected in parallel to the source as shown in figure. Determine the value of $R$ so that maximum power get dissipated in resistance $R$.
\qfig{p1_phy_s1_q07_1}
\begin{qoptions}
\opt{A}{$\dfrac{50}{11}$}
\opt{B}{$\dfrac{50}{12}$}
\opt{C}{50}
\opt{D}{$\dfrac{50}{13}$}
\end{qoptions}
\qn{8} A solid conducting sphere of center `O', radius $2R$, carrying charge $Q$ is surrounded by two point charges $Q$ and $2Q$ as shown in the figure. The electric field at point $P$ due to the total charges of conducting sphere is: ($OP=R$ and $K=\dfrac1{4\pi\epsilon_0}$ )
\qfig{p1_phy_s1_q08_1}
\begin{qoptions}
\opt{A}{$\dfrac7{16}\dfrac{KQ}{R^2}$ towards right}
\opt{B}{$\dfrac18\dfrac{KQ}{R^2}$ towards right}
\opt{C}{$\dfrac{KQ}{R^2}$ towards right}
\opt{D}{zero}
\end{qoptions}
\qn{9} Potential energy curve (sinusoidal) is shown graphically for a particle. The potential energy does not depend on $y$ and $z$ coordinates. The maximum value of corresponding conservative force (in magnitude) for range $0<x<2$ is $(\beta\pi)$. Find the value of $\beta$.
\qfig{p1_phy_s1_q09_1}
\begin{qoptions}
\opt{A}{3N}
\opt{B}{6N}
\opt{C}{4N}
\opt{D}{8N}
\end{qoptions}
\qn{10} An insulating cylindrical rod of diameter $d$ and length $\ell$ ($\ell\gg d$) has a uniform surface charge density such that the electric field just outside the curved surface of the cylinder at point $M$ is $E_0$. Find the electric field due to charge distribution at point $P$ ($r\gg\ell$).
\qfig{p1_phy_s1_q10_1}
\begin{qoptions}
\opt{A}{$E_0\dfrac{\ell d}{2r^2}$}
\opt{B}{$E_0\dfrac{\ell d}{4r^2}$}
\opt{C}{$E_0\dfrac{\ell d}{3r^2}$}
\opt{D}{$E_0\dfrac{2\ell d}{r^2}$}
\end{qoptions}
\sect{SEC--2}
\qn{1} A rod $CD$ of length $L$ and mass $M$ is placed horizontally on a frictionless horizontal surface as shown. A second identical rod $AB$ which is also placed horizontally (perpendicular to $CD$) on the same horizontal surface is moving along the surface with a velocity $v$ in a direction perpendicular to rod $CD$ and its end $B$ strikes the rod $CD$ at end $C$ and sticks to it rigidly. Then:
\qfig{p1_phy_s2_q01_1}
\begin{qoptions}
\opt{A}{speed of centre of mass of the system just after impact is $\dfrac v4$.}
\opt{B}{the $\omega$ (angular speed) of system just after collision is $\dfrac{3v}{5L}$.}
\opt{C}{speed of centre of mass of the system just after impact is $\dfrac v2$.}
\opt{D}{the $\omega$ (angular speed) of system just after collision is $\dfrac{5v}{3L}$.}
\end{qoptions}
\qn{2} Two adiabatic process involving an ideal gas are plotted on a $P$ - $V$ diagram. $A$ and $B$ are two points on these curves as shown in the diagram. Which of the following statements is/are true?
\qfig{p1_phy_s2_q02_1}
\begin{qoptions}
\opt{A}{Heat is given to system in the process $AB$.}
\opt{B}{Heat is rejected by system in the process $AB$.}
\opt{C}{Whether heat is rejected or absorbed by the system during the process $AB$ does not depend on the magnitude of the slope of line $AB$.}
\opt{D}{Whether heat is rejected or absorbed by the system during the process $AB$ does not depend on the position of the points $A$ and $B$ on the $P$ - $V$ diagram.}
\end{qoptions}
\qn{3} Choose the correct option(s):
\begin{qoptions}
\opt{A}{Waves from two coherent sources may be in same phase always at some points.}
\opt{B}{Waves from two incoherent sources have fixed phase difference at a point always}
\opt{C}{Initial phase difference between the waves emitted by two coherent sources may vary with time.}
\opt{D}{Waves from two coherent sources may be in opposite phase always at some points.}
\end{qoptions}
\qn{4} In the given a.c. circuit, choose the correct statement(s).
\qfig{p1_phy_s2_q04_1}
\begin{qoptions}
\opt{A}{Impedance of circuit is $2\Omega$}
\opt{B}{Power factor of circuit is $\dfrac1{\sqrt2}$}
\opt{C}{Peak value of current through resistance is $50\sqrt2$ A}
\opt{D}{Average power supplied by source is 2500W.}
\end{qoptions}
\qn{5} In the figure shown, $C_1=11\mu$F and $C_2=5\mu$F, then at steady state:
\qfig{p1_phy_s2_q05_1}
\begin{qoptions}
\opt{A}{the potential difference across $C_1$ is 5V.}
\opt{B}{the potential difference across $C_2$ is 2V.}
\opt{C}{the potential difference between points $a$ and $b$ is 4V.}
\opt{D}{the potential difference between the terminals of 15 V battery is 9V.}
\end{qoptions}
\sect{SEC--3}
\qn{1} A stone projected at an angle of $60^\circ$ from the ground level strikes at an angle of $30^\circ$ on the roof of a building of height `$h$'. If the speed of projection of the stone is $\sqrt{Ngh}$, find $N$.
\qfig{p1_phy_s3_q01_1}
\qn{2} A particle is projected from point $A$, that is at a distance $4R$ from the centre of the earth, with speed $v_1$ in a direction making $30^\circ$ with the line joining the centre of the earth and point $A$, as shown. Find $v_1$ if the particle passes grazing the surface of the earth. Consider gravitational interaction only between these two. Express you answer in the form $\dfrac{1000X}{\sqrt2}$ m/s and fill value of $X$. (use $\dfrac{GM}R=6.4\times10^7$ m$^2$/s$^2$)
\qfig{p1_phy_s3_q02_1}
\qn{3} Two tuning forks $A$ and $B$ lying on opposite sides of observer `$O$' and of natural frequency 85Hz move with velocity 10 m/s relative to stationary observer $O$. Fork $A$ moves away from the observer while the fork $B$ move towards him. A wind with a speed 10 m/s is blowing towards fork $A$. Find the beat frequency measured by the observer in Hz. [Take speed of sound in air as 340 m/s]
\qn{4} A ray of light travelling in glass ($\mu=3/2$) is incident on a horizontal glass air surface at the critical angle $\theta_C$. If a thin layer of water ($\mu=4/3$) is now poured on the glass air surface, the ray of light emerge into air at the water air surface at an angle of $\pi/k$ radians from the normal at that point, find the value of $k$.
\qfig{p1_phy_s3_q04_1}
\qn{5} A solid has a volume $v_0$ when external pressure is $p_0$. Bulk modulus of the material is $B$. Now the pressure is increased to $p$. The additional energy per unit volume now stored in the material is $\dfrac{3\left(p^2-p_0^2\right)}{zB}$. Find $z=\,$?
\partsec{2}{CHEMISTRY}
\sect{SEC--1}
\qn{1} The number of waves made by a Bohr electron in an orbit of maximum magnetic quantum number $+2$ is
\begin{qoptions}
\opt{A}{3}
\opt{B}{4}
\opt{C}{2}
\opt{D}{1}
\end{qoptions}
\qn{2} Under which of the following conditions a real gas resembles an ideal gas $-$
\begin{qoptions}
\opt{A}{`$a$' and `$b$' are very large}
\opt{B}{`$a$' and `$b$' are very small}
\opt{C}{`$a$' is large and `$b$' is small}
\opt{D}{`$a$' is small and `$b$' is large}
\end{qoptions}
\qn{3} If valence factor ($n$-factor) of compound $\mathrm{NaHC_2O_4.2H_2C_2O_4.3K_2C_2O_4.4Al_2(C_2O_4)_3.3FeC_2O_4}$ in acid base titration is $x$ and redox titration with $\mathrm{KMnO_4}$ is $y$ then value of $y/x$ is $-$
\begin{qoptions}
\opt{A}{8.4}
\opt{B}{9}
\opt{C}{11.25}
\opt{D}{12}
\end{qoptions}
\qn{4} An ammonia-ammonium chloride buffer has a pH value of 9 with $[\mathrm{NH_3}]=0.25$M. What will be the new pH if 500 ml 0.1 M KOH is added to 200 ml buffer solution ($K_b=2\times10^{-5}$)\\
{[}log2 = 0.3]
\begin{qoptions}
\opt{A}{8.4}
\opt{B}{9.6}
\opt{C}{5.6}
\opt{D}{4.4}
\end{qoptions}
\qn{5}
\qfig{p1_chem_s1_q05_1}
\begin{qoptions}
\opt{A}{\qfigi{p1_chem_s1_q05_A}}
\opt{B}{\qfigi{p1_chem_s1_q05_B}}
\opt{C}{\qfigi{p1_chem_s1_q05_C}}
\opt{D}{\qfigi{p1_chem_s1_q05_D}}
\end{qoptions}
\qn{6}
\qfig{p1_chem_s1_q06_1}
The major product in the given reaction is
\begin{qoptions}
\opt{A}{\qfigi{p1_chem_s1_q06_A}}
\opt{B}{\qfigi{p1_chem_s1_q06_B}}
\opt{C}{\qfigi{p1_chem_s1_q06_C}}
\opt{D}{\qfigi{p1_chem_s1_q06_D}}
\end{qoptions}
\qn{7}
\qfig{p1_chem_s1_q07_1}
$P_1$(Ketone) $+$ $P_2$(Alcohol) True about product of above reaction
\begin{qoptions}
\opt{A}{$P_1$ will give $+$ ve iodoform test}
\opt{B}{Product $P_2$ will give $\mathrm{CO_2}$ on heating}
\opt{C}{Product $P_1$ will give only one oxime with $\mathrm{NH_2OH}$}
\opt{D}{Product $P_1$ will give alcohol with $\mathrm{CH_3MgBr}$ followed by hydrolysis.}
\end{qoptions}
\qn{8} Generally diamonds are good conductors of heat. This property of diamonds can be utilized
\begin{qoptions}
\opt{A}{To estimate the hardness of diamonds}
\opt{B}{To determine the refractive index of diamonds}
\opt{C}{To find out whether the given diamond is fake or original}
\opt{D}{To determine its chemical reactivity}
\end{qoptions}
\qn{9} Which of the following will form $2^\circ$ carbocation?
\begin{qoptions}
\opt{A}{\qfigi{p1_chem_s1_q09_A}}
\opt{B}{\qfigi{p1_chem_s1_q09_B}}
\opt{C}{\qfigi{p1_chem_s1_q09_C}}
\opt{D}{\qfigi{p1_chem_s1_q09_D}}
\end{qoptions}
\qn{10} Synthetic detergents are more effective than ordinary soaps in hard water because
\begin{qoptions}
\opt{A}{They are highly soluble in water}
\opt{B}{Their calcium and magnesium salts are water soluble}
\opt{C}{Their calcium and magnesium salts are insoluble in water}
\opt{D}{They convert calcium and magnesium ions by converting them to stable Complex compounds}
\end{qoptions}
\sect{SEC--2}
\qn{1} Which of the following is/are characteristic of chemisorptions
\begin{qoptions}
\opt{A}{It is irreversible}
\opt{B}{It is specific}
\opt{C}{It is multilayer phenomenon}
\opt{D}{Magnitude of heat of adsorption is 80-240 kJ}
\end{qoptions}
\qn{2} Find compounds given below, which produce metallic residue on heating
\begin{qoptions}
\opt{A}{$\mathrm{Ag_2CO_3}$}
\opt{B}{HgO}
\opt{C}{$\mathrm{BaO_2}$}
\opt{D}{$\mathrm{Ba(N_3)_2}$}
\end{qoptions}
\qn{3} Choose the incorrect option(s)
\begin{qoptions}
\opt{A}{\qfigi{p1_chem_s2_q03_A}}
\opt{B}{\qfigi{p1_chem_s2_q03_B} major product exhibits NaH eq enantiomerism}
\opt{C}{Me $-$ Cl $\xrightarrow[\mathrm{Ag_2O}]{\mathrm{Moist}}$ Me $-$ O $-$ Me (major)}
\opt{D}{MeI (excess) $+$ $\mathrm{NH_3}\to\mathrm{Me_3N}$}
\end{qoptions}
\qn{4} Select correct for hybridization
\begin{qoptions}
\opt{A}{As \% s-character increases in hybrid orbital, bond angle increases}
\opt{B}{As \% s-character decreases in hybrid orbital, bond strength increases}
\opt{C}{As \% s-character increases in hybrid orbital electro negativity increases}
\opt{D}{As \% s-character increases in hybrid orbital, orbital length is decreases}
\end{qoptions}
\qn{5} Which of the following statements are correct?
\begin{qoptions}
\opt{A}{$\mathrm{[Ag(NH_3)_2]^+}$ is less stable than $\mathrm{[Ag(CN)_2]^-}$}
\opt{B}{$\mathrm{[Ag(NH_3)_2]^+}$ is more stable than $\mathrm{[Ag(en)]^+}$}
\opt{C}{AgCl is soluble in conc HCl}
\opt{D}{AgCl converts into AgI when it is dissolve in $\mathrm{NH_4OH}$ followed by the addition of KI(aq.).}
\end{qoptions}
\sect{SEC--3}
\qn{1} Total number of layers between a distance of $14R\sqrt{\dfrac23}$ in hcp unit cell is (including first and last).
\qn{2} Quinhydrone is a sparingly soluble one-to-one addition compound formed from hydroquinone and quinone. When solid quinhydrone is dissolved in an aqueous medium, equal concentrations of hydroquinone and quinone result. If a platinum wire is dipped into the solution, the potential of the electrode is governed by the reversible reaction :
\qfig{p1_chem_s3_q02_1}
Where $Q$ denotes quinone and $\mathrm{H_2Q}$ denotes hydroquinone. The $E^\circ$ value of this half cell reaction is $+0.46$V with respect to the saturated calomel reference electrode (SCE). Because the potential of the quinine-hydroquinone half reaction depends on the concentration of hydrogen ion, it is possible to use this system to measure pH. A sample solution of unknown pH was saturated with quinohydrone and a platinum electrode was dipped into it. If the potential of the such electrode was found to be $+0.22$V with respect to SCE, what was the PH of the sample solution?\\
{[}Given : $\dfrac{2.303\,RT}{F}=0.06$]
\qn{3}
\qfig{p1_chem_s3_q03_1} number of iodoform molecules produced per molecule.
\qn{4} The electro negativity difference between two elements $A$ and $B$ is 0.5. the magnitude of percent ionic character in single $A$--$B$ would be nearest to
\qn{5} $\mathrm{PCl_5}$ reacts with $\mathrm{P_4O_{10}}$ to give a single product. The number of $\mathrm{PCl_5}$ molecules required to react with one mole of $\mathrm{P_4O_{10}}$ in the balanced reaction is.
\partsec{3}{MATHEMATICS}
\sect{SEC--1}
\qn{1} The term independent of $x$ and $y$ in the expansion of $\Bigl[\left(\sqrt x+\dfrac1{\sqrt x}\right)^2+\allowbreak\left(\sqrt y+\dfrac1{\sqrt y}\right)^2+\allowbreak\left(\sqrt{xy}+\dfrac1{\sqrt{xy}}\right)^2-4\Bigr]^n$ is
\begin{qoptions}
\opt{A}{$\left(\displaystyle\sum_{r=0}^{n}{}^nC_r\right)^2$}
\opt{B}{$\displaystyle\sum_{r=0}^{n}\left({}^nC_r\right)^2$}
\opt{C}{$\left(\displaystyle\sum_{r=0}^{n}{}^nC_r\right)^3$}
\opt{D}{$\displaystyle\sum_{r=0}^{n}\left({}^nC_r\right)^3$}
\end{qoptions}
\qn{2} If $y=2\cot^{-1}\left(\dfrac x2\right)+x$ then $\displaystyle\int\dfrac{x^2\,dx}{(x^2-4)\sin x+4x\cos x}=$\\
(Where $C$ is constant of integration)
\begin{qoptions}
\opt{A}{$\ln|\mathrm{cosec}2y+\cot y|+C$}
\opt{B}{$\ln|\mathrm{cosec}2y-\cot y|+C$}
\opt{C}{$\ln|\mathrm{cosec}y+\cot y|+C$}
\opt{D}{$\ln|\mathrm{cosec}y-\cot y|+C$}
\end{qoptions}
\qn{3} Consider a curve $|x|+|y|=1$ such that lines $y=mx$, $y=nx$ make points of intersection in the same quadrant. Let $A$, $B$ be two such points lying in the same quadrant such that $OA$, $AB$, $OB$ are in A.P. If $\theta$ is angle between $OA$ and $OB$ then the maximum angle $\theta$ is ($O$ is origin)
\begin{qoptions}
\opt{A}{$\dfrac\pi6$}
\opt{B}{$\dfrac\pi3$}
\opt{C}{$\dfrac\pi4$}
\opt{D}{None of these}
\end{qoptions}
\qn{4} $n$ balls are arranged to form an equilateral triangle such that first row consists of 1 ball, second row consist of 2 balls, third row consists of 3 balls and so on 49 more balls are added to form a square. If each side of square contains 3 balls more the number of balls each side of the equilateral triangle contain then $n=$
\begin{qoptions}
\opt{A}{10}
\opt{B}{20}
\opt{C}{50}
\opt{D}{15}
\end{qoptions}
\qn{5} A line $L$ cuts the lines $AB$, $AC$ and $AD$ of a parallelogram $ABCD$ at points $B_1$, $C_1$ and $D_1$ respectively. If $\overrightarrow{AB_1}=\lambda_1\overrightarrow{AB}$, $\overrightarrow{AD_1}=\lambda_2\overrightarrow{AD}$, $\overrightarrow{AC_1}=\lambda_3\overrightarrow{AC}$ then $\dfrac1{\lambda_3}=$
\begin{qoptions}
\opt{A}{$\dfrac1{\lambda_1}+\dfrac1{\lambda_2}$}
\opt{B}{$\dfrac1{\lambda_1}-\dfrac1{\lambda_2}$}
\opt{C}{$-\lambda_1+\lambda_2$}
\opt{D}{$\lambda_1+\lambda_2$}
\end{qoptions}
\qn{6} If $[x^3+x^2+x+1]=[x^3+x^2+1]+x$, (where $[.]$ denotes the greatest integer function), then number of solutions of the equation $\log|[x]|=2-|[x]|$ is
\begin{qoptions}
\opt{A}{1}
\opt{B}{0}
\opt{C}{3}
\opt{D}{2}
\end{qoptions}
\qn{7} Let the circle $(x-1)^2+(y-2)^2=25$ cuts rectangular hyperbola with transverse axis along $y=x$ at 4 points $A,B,C$ and $D$ having coordinates $(x_i,y_i)$, $i=1,2,3,4$ such that origin is centre of hyperbola. If $\ell=x_1+x_2+x_3+x_4$, $m=x_1^2+x_2^2+x_3^2+x_4^2$, $n=y_1^2+y_2^2+y_3^2+y_4^2$ then $\dfrac{m+n}\ell=$
\begin{qoptions}
\opt{A}{100}
\opt{B}{10}
\opt{C}{50}
\opt{D}{20}
\end{qoptions}
\qn{8} If $P=\left(\tan3^{n+1}\theta-\tan\theta\right)$ and $Q=\displaystyle\sum_{r=0}^{n}\dfrac{\sin3^r\theta}{\cos3^{r+1}\theta}$ then
\begin{qoptions}
\opt{A}{$P=3Q$}
\opt{B}{$P=2Q$}
\opt{C}{$2P=Q$}
\opt{D}{$3P=Q$}
\end{qoptions}
\qn{9} $A,B,C$ are vertices of a triangle with right angle at $A$ and $P(-4,0)$ ; $Q(0,6)$ are two given points. If the ratio of distances from each vertex to $P$, to that of $Q$ is $2:3$, then the centroid of $\Delta ABC$ lies on a circle with radius equal to
\begin{qoptions}
\opt{A}{$\dfrac{4\sqrt{13}}5$ units}
\opt{B}{4 units}
\opt{C}{$\dfrac{8\sqrt{13}}5$ units}
\opt{D}{8 units}
\end{qoptions}
\qn{10} A line is drawn from the point $P(1,1,1)$ and perpendicular to a line with direction ratios $(1,1,1)$ to intersect the plane $x+2y+3z=4$ at $Q$. The locus of point $Q$ is
\begin{qoptions}
\opt{A}{$\dfrac x1=\dfrac{y-5}{-2}=\dfrac{z+2}1$}
\opt{B}{$\dfrac x{-2}=\dfrac{y-5}1=\dfrac{z+2}1$}
\opt{C}{$x=y=z$}
\opt{D}{$\dfrac x2=\dfrac y3=\dfrac z5$}
\end{qoptions}
\sect{SEC--2}
\qn{1} Consider $\Delta ABC$ whose circumcircle is $|z|=r$ in the argand plane with $A,B,C$ be represented by complex numbers $a,b,c$ respectively. The foot of altitudes from $A,B,C$ meet the opposite side at $D,E$ and $F$ and the altitudes when produced meet the circle $|z|=r$ at $L,M$ and $N$ respectively then which of following is/are correct ?
\begin{qoptions}
\opt{A}{centroid of $\Delta LMN$ is $\dfrac{-\left(a^2b^2+b^2c^2+c^2a^2\right)}{3abc}$}
\opt{B}{the mirror image of orthocenter of $\Delta ABC$ with respect to $BC$ lies on $|z|=r$}
\opt{C}{the mirror image of orthocenter of $\Delta ABC$ with respect to $BC$ is $\dfrac{-bc}a$}
\opt{D}{none of these}
\end{qoptions}
\qn{2} If $f(x)+f(y)=\dfrac1x+\dfrac1y$, $\forall\ x,y\in\mathbb R-\{0\}$ and $\displaystyle\int_2^3\dfrac{3(f(x))^5-f(x)}{1-(f(x))^4}\,dx=\dfrac12\log\dfrac{2^\alpha}{3^\beta}$, then
\begin{qoptions}
\opt{A}{$\alpha>\beta$}
\opt{B}{$\beta$ is prime}
\opt{C}{$\alpha<2\beta$}
\opt{D}{$(\alpha+\beta)$ is prime}
\end{qoptions}
\qn{3} Three students $A,B,C$ appear for IIT-JEE exam. Their respective probabilities of clearing the exams being $\dfrac3{10}$, $\dfrac12$ and $\dfrac25$. Each student clearing or not clearing the exam is independent of other. If the probability of $A$ passing the exam be $\dfrac mn$ ($m,n\in N$) given that only one of the student has cleared the exam, then the least value of $(m+n)$ is greater than
\begin{qoptions}
\opt{A}{10}
\opt{B}{20}
\opt{C}{40}
\opt{D}{80}
\end{qoptions}
\qn{4} For a given parabola $y^2=4ax$, two variable chords $PQ$ and $RS$ at right angles are drawn through the fixed point $A(x_1,y_1)$ inside the parabola, making variable angles $\theta$ and $\alpha$ with $x$-axis. If $r_1,r_2,r_3,r_4$ are distances of $P,Q,R$ and $S$ from $A$, then the value of $\dfrac1{r_1r_2}+\dfrac1{r_3r_4}$
\begin{qoptions}
\opt{A}{independent of $\theta$}
\opt{B}{independent of $\alpha$}
\opt{C}{depends upon both $\theta$ \& $\alpha$}
\opt{D}{is a constant}
\end{qoptions}
\qn{5} Let $f(x)=|(x-1)|+|(x-4)|+|(x-9)|+\ldots..+|(x-2401)|+|(x-2500)|\ \forall\ x\in\mathbb R$. If $m$ and $n$ denote the number of points of non differentiability of $f(x)$ and number of integral points where $f(x)$ has minimum value respectively then
\begin{qoptions}
\opt{A}{$m=50$}
\opt{B}{$n=52$}
\opt{C}{$m+n=100$}
\opt{D}{$m=n=50$}
\end{qoptions}
\sect{SEC--3}
\qn{1} Let $x,y\in\mathbb R$ such that $\cos x\cos y+2\sin y+2\sin x\cos y=3$, then find the value of $\tan^2x+5\tan^2y$.
\qn{2} Suppose $f:\mathbb R\to\mathbb R$ be a continuous function such that $f(3x)=f(x)$, $\forall\ x\in\mathbb R$ and $f(2)=5$, then value of $\displaystyle\int_2^5\operatorname{sgn}(f(x))\,dx$ is equal to (where $\operatorname{sgn}(x)$ represents signum function)
\qn{3} If $\log_{245}175=a$, $\log_{1715}875=b$, then the value of $\dfrac{1-ab}{a-b}$ is
\qn{4} The value of the expression
\begin{align*}&\lim_{y\to\infty}\Bigl((8y+3)(y+1)(y+2)(y+3)(2y+3)\\&\quad(4y+5)(2y+7)\Bigr)^{1/7}-2y\end{align*}
is of form $\dfrac AB$ (Where $A$, $B$ are coprime number). Then the value $\left[\dfrac AB\right]$ (where $[.]$ denotes the greatest integer function), is
\qn{5} The plane $\dfrac x1+\dfrac y2+\dfrac z3=1$ intersect $x$-axis, $y$-axis, $z$-axis at $A,B,C$ respectively. If the distance between origin and othrocenter of $\Delta ABC$ is equal to $k$ then value of $7k$ is equal to
\papersec{2}
\partsec{1}{PHYSICS}
\sect{SEC--1}
\qn{1} A solid uniform ball of mass $m$ and radius $R$ collides elastically with a rough fixed inclined surface as shown. Velocity and angular velocity of ball just before collision are $V_0$ and $\omega_0=\dfrac{5V_0}R$. Choose correct option(s) -
\qfig{p2_phy_s1_q01_1}
\begin{qoptions}
\opt{A}{If $\mu=\dfrac34$ then the velocity of ball just after collision is vertically upwards.}
\opt{B}{If $\mu=\dfrac34$ then the velocity of ball just after collision is horizontal leftwards.}
\opt{C}{If $\mu=\dfrac34$ then angular velocity after collision is $\dfrac{2V_0}R$}
\opt{D}{If $\mu=\dfrac34$ then velocity of ball after collision is $V_0$.}
\end{qoptions}
\qn{2} The resistance per unit length of potentiometer wire of uniform cross section is $\left(\dfrac{3R_0x}{\ell^2}\right)$, where $x$ is measured from end $A$. Balance length is also measured from $A$. Select correct option (s) :
\qfig{p2_phy_s1_q02_1}
\begin{qoptions}
\opt{A}{When $S$ is open balance length is at $\dfrac\ell2$.}
\opt{B}{When $S$ is open balance length is at $\ell\sqrt{\dfrac56}$.}
\opt{C}{When $S$ is closed balance length is at $\ell\sqrt{\dfrac34}$.}
\opt{D}{When $S$ is closed balance length is at $\ell\sqrt{\dfrac78}$.}
\end{qoptions}
\qn{3} The pitch of screw gauge is 1 mm and its circular scale is divided into 100 divisions. When nothing is put between the studs the zero of main scale is not seen but when circular scale is rotated by $450^\circ$ the zero of main scale is just visible and the zero of main scale coincides with the zero of circular scale. When a glass plate is placed between the studs, the circular scale lies between $18^{\text{th}}$ and $19^{\text{th}}$ division of main scale and circular scale reads 34 divisions. Then,
\begin{qoptions}
\opt{A}{There is positive zero error and its magnitude is 1.25 mm.}
\opt{B}{There is negative zero error and its magnitude is 1.25 mm.}
\opt{C}{The thickness of the glass plate is 19.59 mm.}
\opt{D}{The thickness of the glass plate is 17.09 mm.}
\end{qoptions}
\qn{4} A semi circular disc of radius $R$ is placed on a smooth horizontal surface. Figure shows top view of this. If temperature is increased by $\Delta t$ then displacement of (coefficient of linear expansion is $\alpha$) :
\qfig{p2_phy_s1_q04_1}
\begin{qoptions}
\opt{A}{$A$ along $y$-axis is $R\alpha\Delta t$.}
\opt{B}{$A$ along $x$-axis is $-\dfrac{4R}{3\pi}\alpha\Delta t$.}
\opt{C}{$B$ along $x$-axis is $R\alpha\Delta t$.}
\opt{D}{$B$ along $y$-axis is zero.}
\end{qoptions}
\sect{SEC--2}
\paraq{1}{2}
Consider the given system shown in figure. At certain instant velocity of lift (w.r.t. ground), velocity of $A$ (relative to lift) is shown. Velocity of lift is constant.
\qfig{p2_phy_s2_para1_2_1}
\qn{1} Choose the correct options.
\begin{qoptions}
\opt{A}{Velocity of $A$ relative to ground is 6 m/s upward.}
\opt{B}{Velocity of $B$ relative to ground is 2 m/s upward.}
\opt{C}{Velocity of $B$ relative to ground is 4 m/s downward.}
\opt{D}{Velocity of $A$ relative to ground is 4 m/s upward.}
\end{qoptions}
\qn{2} If $m_A=m_B=m$, choose the correct options.
\begin{qoptions}
\opt{A}{acceleration of $A$ is $\dfrac g3$ downward.}
\opt{B}{acceleration of $A$ is $\dfrac g5$ upward.}
\opt{C}{force on the ceiling of lift by string is $\dfrac{6mg}5$.}
\opt{D}{acceleration of $B$ is $\dfrac g5$ downward.}
\end{qoptions}
\paraq{3}{4}
An infinitely long wire lying along $z$-axis carries a current $I$, flowing towards positive $z$-direction. There is no other current, consider a circle in $x$-$y$ plane with centre at (2 meter, 0, 0) and radius 1 meter. Divide the circle in small segments and let $d\vec\ell$ denote the length of a small segment in anticlockwise direction, as shown.
\qfig{p2_phy_s2_para3_4_1}
\qn{3} Consider two points $A(3,0,0)$ and $B(2,1,0)$ on the given circle. The path integral $\displaystyle\int_A^B\vec B\cdot d\vec\ell$ of the total magnetic field $\vec B$ along the perimeter of the given circle from $A$ to $B$ is,
\begin{qoptions}
\opt{A}{$\dfrac{\mu_0I}\pi\tan^{-1}\dfrac12$}
\opt{B}{$\dfrac{\mu_0I}{2\pi}\tan^{-1}\dfrac12$}
\opt{C}{$\dfrac{\mu_0I}{2\pi}\sin^{-1}\dfrac12$}
\opt{D}{0}
\end{qoptions}
\qn{4} The maximum value of path integral $\displaystyle\int\vec B\cdot d\vec\ell$ of the total magnetic field $\vec B$ along the perimeter of the given circle between any two points on the circle is
\begin{qoptions}
\opt{A}{$\dfrac{\mu_0I}{12}$}
\opt{B}{$\dfrac{\mu_0I}8$}
\opt{C}{$\dfrac{\mu_0I}6$}
\opt{D}{0}
\end{qoptions}
\paraq{5}{6}
Resistance per unit length of wire frame $ABCDA$ is $\dfrac R{2\ell}$. The wire frame is moving with constant velocity $V_0$ on two parallel conducting and smooth horizontal rails in uniform magnetic field $B$ perpendicular to plane of rails as shown in figure.
\qfig{p2_phy_s2_para5_6_1}
\qn{5} Select correct option :-
\begin{qoptions}
\opt{A}{Current in resistor $R$ is $\dfrac{2BV_0\ell}{5R}$.}
\opt{B}{Current in side $AB$ is $\dfrac{4BV_0\ell}{5R}$.}
\opt{C}{Current in side $CD$ is $\dfrac{4BV_0\ell}{5R}$.}
\opt{D}{Current in side $ST$ is $\dfrac{2BV_0\ell}{5R}$.}
\end{qoptions}
\qn{6} Select correct option :-
\begin{qoptions}
\opt{A}{External force required to move to frame is $\dfrac{2B^2\ell^2V_0}{5R}$.}
\opt{B}{External force required to move to frame is $\dfrac{4B^2\ell^2V_0}{5R}$.}
\opt{C}{Potential difference between $B$ and $C$ is zero.}
\opt{D}{Potential difference between $B$ and $C$ is $2BV_0\ell$.}
\end{qoptions}
\sect{SEC--3}
\qn{1} A Physical quantity $P$ is given as - $P=\dfrac{x^2y}{(4-t)^3}$ where, $x=(2\pm0.01)$ SI unit; $y=(4\pm0.02)$ SI unit; $t=(2\pm0.01)$ SI unit. Find the percentage error in calculation of $P$.
\qn{2} A current carrying string in shape in ring of radius 1m is placed in uniform magnetic field perpendicular to plane of ring. The current and area of cross-section of ring are 1 amp and 0.2cm$^2$. The density of string is $2\times10^3$ kg/m$^3$. The speed of wave in the string (ring) is 10 m/sec. Find magnetic field (in T)? (Neglect any other interaction)
\qfig{p2_phy_s3_q02_1}
\qn{3} Initially $S$ is open and capacitor $A$ is charged up to potential difference $\dfrac V2$ and capacitor $B$ is uncharged. After closing the switch $S$ the heat loss in the circuit is $\dfrac{CV^2}{2\times n}$J. Find $n$?
\qfig{p2_phy_s3_q03_1}
\qn{4} Two parallel square plates each of side length `$a$' are arranged such that line joining centre of both the plates is perpendicular to planes of both plates. Eight fixed charges are placed at corners of plate as shown in figure. A charge particle of charge $+q$ and mass $m$ is projected with initial velocity $v_0$ from centre $C_1$ of one square plate towards centre $C_2$ of other. The length of line $C_1C_2$ is $4a$. The speed of charge particle at mid-point of line $C_1C_2$ is $\dfrac{4v_0}x$. Find the value of $x$ ? (Neglect effect of gravity).
\qfig{p2_phy_s3_q04_1}
\qn{5} The density of a solid spherical planet of radius $R$ is given as $\rho=\rho_0r$, where $\rho_0=$ constant and $r$ is distance measured from centre of planet. The acceleration due to gravity of this planet is half of maximum value at distance $x$ from centre and also at a distance $y$ from the centre. The value of $x+y$ is $R\sqrt{\dfrac\alpha\beta}$. Here $\alpha$ and $\beta$ are single digit integer. Find the value of $\alpha$-$\beta$?
\qfig{p2_phy_s3_q05_1}
\qn{6} A uniform vertical cylinder (density $=\sigma$) is released from rest when its lower end just touches the liquid (density $=\rho$) surface of a deep lake. Calculate maximum displacement of cylinder (in meter) Take, $\ell=8$m and $\dfrac\sigma\rho=\dfrac12$
\qfig{p2_phy_s3_q06_1}
\qn{7} A smooth cylindrical glass lies on the ground. A small particle of mass $m$ lies on the bottom surface touching circumference as shown in the figure. The minimum value of $V_0$ (along the circumference) such that particle just reaches at the top of glass is $\sqrt{\dfrac{Kgh}3}$. Find the value of $K$.
\qfig{p2_phy_s3_q07_1}
\qn{8} The square of speed of a particle moving along straight line varying with time as shown in graph. The equation of graph is $\upsilon=\dfrac{C^2}{\sqrt t}$, where $C=$ constant, ($\upsilon$ in meter/sec. and $t$ in sec.) Find the value of constant $C$ in SI units.
\qfig{p2_phy_s3_q08_1}
\qn{9} Two radioactive samples $A_1$ and $A_2$ having half life 3 years and 2 years respectively have been decaying for many years. Today the number of atoms in the sample $A_1$ is twice the number of atoms in the sample $A_2$. Both the samples had same number of atoms $X$ years ago, calculate $X$.
\qn{10} Two conductors $A$ and $B$ each of cross section area 5 cm$^2$ are connected in series. Variation of temperature (in $^\circ$C) along the length (in cm) is as shown in the figure. If thermal conductivity of $A$ is 120 J/m-sec-$^\circ$C. The thermal conductivity of $B$ is $40K$ J/m-sec-$^\circ$C. Find the value of $K$.
\qfig{p2_phy_s3_q10_1}
\partsec{2}{CHEMISTRY}
\sect{SEC--1}
\qn{1} $\lambda$ of 0.01 M $\mathrm{MgCl_2}$ is 220 S.cm$^2$.mol$^{-1}$\\
$\lambda^\circ$ of $\mathrm{MgCl_2}$ is 240 S.cm$^2$.mol$^{-1}$\\
$\lambda^\circ$ of $\mathrm{Mg^{2+}}$ is 100 S.cm$^2$.mol$^{-1}$\\
All above measurements are done in an conductivity cell.\\
Assuming $\mathrm{MgCl_2}$ to be a strong electrolyte, which of the following is incorrect?
\begin{qoptions}
\opt{A}{Limiting equivalent conductivity of $\mathrm{MgCl_2}$ is 480 S.cm$^2$.eq$^{-1}$}
\opt{B}{$\lambda$ of $\mathrm{MgCl_2}$ varies linearly with its concentration}
\opt{C}{$\lambda$ of 0.04 M $\mathrm{MgCl_2}$ will be 160 S.cm$^2$.mol$^{-1}$}
\opt{D}{$\lambda$ of 0.04 M $\mathrm{MgCl_2}$ will be 200 S.cm$^2$.mol$^{-1}$}
\end{qoptions}
\qn{2} When two phases of the same single substance remain in equilibrium with one another at a constant P and T, their molar------must be equal. Which of the following will not fit into the blank ?
\begin{qoptions}
\opt{A}{Internal energy}
\opt{B}{Enthalpy}
\opt{C}{Entropy}
\opt{D}{Gibb's free energy}
\end{qoptions}
\qn{3} Which of the following is/are correct name of the given polymer -
\begin{qoptions}
\opt{A}{$\left[\mathrm{CH_2-\overset{CH_3}{\underset{|}{CH}}}\right]_n$ \quad Polypropene}
\opt{B}{$\left[\mathrm{CH_2-\overset{C_6H_5}{\underset{|}{CH}}}\right]_n$ \quad Polystyrene}
\opt{C}{$\left[\mathrm{NH-CO-NH-CH_2}\right]_n$ \quad Urea-formaldehyde Resin}
\opt{D}{\qfigi{p2_chem_s1_q03_D} \quad Bakelite}
\end{qoptions}
\qn{4} Select the correct statements(s)-
\begin{qoptions}
\opt{A}{\qfigi{p2_chem_s1_q04_A}\\ is a dipeptide of Isoleucine \& Alanine, whose abbreviated name ILE-ALA}
\opt{B}{Compound \qfigi{p2_chem_s1_q04_2} show keto-enol tautomerism}
\opt{C}{Phenol \& benzoic acid can be distinguished by $\mathrm{NaHCO_3}$}
\opt{D}{Order of basicity in aqueous medium $\mathrm{MeNH_2<Me_2NH<Me_3N}$}
\end{qoptions}
\sect{SEC--2}
\paraq{1}{2}
Consider the following redox reaction :\\
$\mathrm{H_2O + AX + BY \longrightarrow HA + YO + BX_2}$ (unbalanced)\\
It is also known that oxidation number of X is $-2$ and neither X nor water is involved in the redox process. (In compound BY, B is cation and consider its oxidation number $<4$).\\
Answer the following questions based on the information given.
\qn{1} Select the incorrect option :-
\begin{qoptions}
\opt{A}{The element A is undergoing reduction}
\opt{B}{The element B is undergoing reduction}
\opt{C}{The element B is undergoing oxidation.}
\opt{D}{The element Y is undergoing oxidation.}
\end{qoptions}
\qn{2} If the given reaction is balanced with smallest whole number coefficients, then-
\begin{qoptions}
\opt{A}{Sum of stoichiometric coefficient of all the components is 8}
\opt{B}{Sum of stoichiometric coefficient of all the components is 9}
\opt{C}{Stoichiometric coefficient of AX is 6}
\opt{D}{Stoichiometric coefficient of HA is 4}
\end{qoptions}
\paraq{3}{4}
Copper is the most noble of the first row transition metals and occurs in small deposits in several countries. Ores of copper include chalcanthite $\mathrm{(CuSO_4.5H_2O)}$, atacamite $\mathrm{(Cu_2Cl(OH)_3)}$, cuprite $\mathrm{(Cu_2O)}$, copper glance $\mathrm{(Cu_2S)}$ and malachite $\mathrm{(Cu_2(OH)_2CO_3)}$. However 80\% of the world copper production comes from the ore chalcopyrites $\mathrm{(CuFeS_2)}$. The extraction of copper from chalcopyrites involves partial roasting, removal of iron and self reduction
\qn{3} Partial roasting of chalcopyrites with silica added to it, in this process does not produces :- .
\begin{qoptions}
\opt{A}{$\mathrm{FeSiO_3}$}
\opt{B}{$\mathrm{FeS}$}
\opt{C}{$\mathrm{Cu_2S}$}
\opt{D}{$\mathrm{CuO}$}
\end{qoptions}
\qn{4} In self reduction the reducing agent is \ldots\ldots\ldots .\\
Select the correct option for blank. -
\begin{qoptions}
\opt{A}{$\mathrm{Cu^+}$}
\opt{B}{$\mathrm{O^{-2}}$}
\opt{C}{$\mathrm{S^{-2}}$}
\opt{D}{$\mathrm{SiO_2}$}
\end{qoptions}
\paraq{5}{6}
An organic compound (X) $\mathrm{C_8H_{10}O}$ was subjected to a series of tests in the laboratory. It was found that this compound
\begin{enumerate}[label=(\roman*),leftmargin=2em,itemsep=0pt,topsep=2pt]
\item Rotates the plane polarized light
\item Evolves hydrogen gas with Na metal
\item Reacts with $\mathrm{I_2}$ and NaOH to produce pale yellow precipitate.
\item Does not react with $\mathrm{Br_2/CCl_4}$
\item Reacts with PCC to form (Y) $\mathrm{C_8H_8O}$, which can also be synthesized by the reaction of benzene and ethanoyl chloride with anhydrous $\mathrm{AlCl_3}$.
\item Loses optical activity when it react with $\mathrm{MnO_2/\Delta}$ followed by Zn-Hg/Conc. HCl to form compound (Z).
\item It gives +ve test with Lucas reagent
\end{enumerate}
\qn{5} The compound X is :
\begin{qoptions}
\opt{A}{\qfigi{p2_chem_s2_q05_A}}
\opt{B}{\qfigi{p2_chem_s2_q05_B}}
\opt{C}{\qfigi{p2_chem_s2_q05_C}}
\opt{D}{\qfigi{p2_chem_s2_q05_D}}
\end{qoptions}
\qn{6} Compound Y and compound Z are respectively -\\ \hspace*{3.2cm}\textbf{Y}\hspace{4cm}\textbf{Z}
\begin{qoptions}
\opt{A}{\qfigi{p2_chem_s2_q06_A}}
\opt{B}{\qfigi{p2_chem_s2_q06_B}}
\opt{C}{\qfigi{p2_chem_s2_q06_C}}
\opt{D}{\qfigi{p2_chem_s2_q06_D}}
\end{qoptions}
\sect{SEC--3}
\qn{1} How many of the following contain peroxide linkage ?\\
(a) $\mathrm{[FeO_4]^{2-}}$ \quad (b) $\mathrm{Mn_2O_7}$ \quad (c) $\mathrm{HSO_5^-}$ \quad (d) $\mathrm{CrO_5}$ \quad (e) $\mathrm{PbO_3^{2-}}$ \quad (f) $\mathrm{HNO_4}$ \quad (g) $\mathrm{Na_2B_4O_7}$ \quad (h) $\mathrm{MnO_2}$ \quad (i) $\mathrm{S_4O_6^{2-}}$
\qn{2} How many of the following orders are correct with respect to their properties indicated ?\\
(i) $\mathrm{HCl<HF}$ \quad Boiling point\\
(ii) $\mathrm{PH_3<NH_3}$ \quad Dipole moment\\
(iii) $\mathrm{H_2O<D_2O}$ \quad Strength of Hydrogen bonding\\
(iv) $\mathrm{F_2<Cl_2}$ \quad Bond energy\\
(v) $\mathrm{H_2O<OCl_2}$ \quad Bond angle\\
(vi) $\mathrm{H_2O_2<O_2F_2}$ \quad Bond length of O--O bond\\
(vii) $\mathrm{NaCl<CuCl}$ \quad Covalent character.\\
(viii) $\mathrm{SO_2<SO_3}$ \quad Acidic strength
\qn{3} Two moles of an ideal monoatomic gas is taken from state A to state B through a process AB in which the pressure and the temperature of the gas vary such that $PT$ = constant. The process can be represented on a P--T graph as follows ($T_1$ = 300 K) :
\qfig{p2_chem_s3_q03_1}
The magnitude of heat involved during process AB is (Give your answer after dividing by 300R) :
\qn{4} The volume of nitrogen gas (at STP) required to cover a sample of silica gel with a monomolecular layer is 112 cm$^3$g$^{-1}$ of gel. Calculate the surface area per gram (in m$^2$) of the gel if each nitrogen molecule occupies $16.0\times10^{-22}$m$^2$. Give your answer to the nearest single digit integer $\{N_A=6\times10^{23}\}$.
\qn{5} Number of crystal system having at least 2 interfacial angles equal is x\\
Number of crystal system having none of the interfacial angles equal is y\\
Number of crystal system having none of the axial length to be equal is z\\
Find the value of $x-y+z$.
\qn{6} What is the value of ``n'' in $\mathrm{Na_2Mg_3Al_n}$\allowbreak$\mathrm{[(Si_4O_{11})_2](OH)_2}$ ?
\qn{7} An alcohol (A) on treatment with conc. $\mathrm{H_2SO_4}$/$\Delta$ gave an alkene (B). The compound (B), on reacting with $\mathrm{Br_2/CCl_4}$ and subsequent didehydrobromination with $\mathrm{NaNH_2}$(2eq.) produced a compound (C). The compound (C) with dil. $\mathrm{H_2SO_4}$ in presence of $\mathrm{HgSO_4}$ gave a compound `D'. The compound D can also be obtained by oxidation of A by pcc or from distillation of calcium acetate.\\
If molecular weight of `C' is $Y\times10$, then write the value of $Y$.
\qn{8} Observe the following reaction\\
$\mathrm{X}\xrightarrow{\mathrm{NH_2OH/HCl}}\mathrm{Y}\xrightarrow{\mathrm{H_2SO_4}}\mathrm{Z}\xrightarrow[\mathrm{(ii)\,H^+}]{\mathrm{(i)\,NaOH/Boil}}\mathrm{H_2N-CH_2-CH_2-CH_2-COOH}$\\
If the molar mass of X is $M\times10$. Then the value of $M$ is:
\qn{9} The total number of all possible Mono Aldols ($\beta$-Hydroxy carbonyl compounds) are formed in the following reaction. (Ignore stereoisomers)
\qfig{p2_chem_s3_q09_1}
\qn{10} The number of stereoisomers of organic products containing $-$C=N--NH$-$ obtained in the following reaction in acidic medium.
\qfig{p2_chem_s3_q10_1}
\partsec{3}{MATHEMATICS}
\sect{SEC--1}
\qn{1} If $\mathrm{f(x)}=\dfrac{[\mathrm x]+1}{\{\mathrm x\}+1}$ for $f:\left[0,\dfrac52\right]\to\left(\dfrac12,3\right]$. Then which of the following is true\\
(where [.] denotes greatest integer function and \{.\} denotes fractional part function)
\begin{qoptions}
\opt{A}{$f(x)$ is injective, discontinuous function}
\opt{B}{$f(x)$ is surjective, non differentiable function}
\opt{C}{$\min\left(\lim\limits_{x\to1^-}f(x),\ \lim\limits_{x\to1^+}f(x)\right)=f(1)$}
\opt{D}{minimum value of $f(x)=\dfrac12$}
\end{qoptions}
\qn{2} A ray of light comes along line L and strikes the plane mirror kept along the plane $P=0$ at B. $A(2,1,6)$ is a point on the line L whose image about $P=0$ is $A'$. It is given that line L is $\dfrac{x-2}3=\dfrac{y-1}4=\dfrac{z-6}5$ and plane $P=0$ is $x+y-2z=3$ then
\begin{qoptions}
\opt{A}{The coordinates of $A'$ are $(6,5,-2)$}
\opt{B}{The coordinates of $A'$ are $(6,5,2)$}
\opt{C}{Equation of the reflected ray is $\dfrac{x+10}4=\dfrac{y+15}5=\dfrac{z+14}3$}
\opt{D}{Equation of the reflected ray is $\dfrac{x+10}4=\dfrac{y-5}4=\dfrac{z+2}3$}
\end{qoptions}
\qn{3} If $\vec a,\vec b,\vec c$ be non-coplanar unit vectors equally inclined to one another at an acute angle $\theta$ if $\vec a\times\vec b+\vec b\times\vec c=p\vec a+q\vec b+r\vec c$, then
\begin{qoptions}
\opt{A}{$p=r$}
\opt{B}{$p=\dfrac1{\sqrt{1+2\cos\theta}},\ q=-\dfrac{2\cos\theta}{\sqrt{1+2\cos\theta}}$}
\opt{C}{$r=\dfrac1{\sqrt{1+2\cos\theta}}$}
\opt{D}{$p=\dfrac{2\cos\theta}{\sqrt{1+2\cos\theta}}$}
\end{qoptions}
\qn{4} From a point $P$ representing complex number $z_1$ on the curve $(|z|=2)$ two tangents are drawn form $P$ to the curve $|z|=1$ meets at $A(z_2)$ and $B(z_3)$ then
\begin{qoptions}
\opt{A}{centroid of $\Delta ABP$ will be on curve $|z|=1$}
\opt{B}{$\arg\left(\dfrac{z_2}{z_3}\right)=\pm\dfrac{2\pi}3$}
\opt{C}{$\left(\dfrac4{\bar z_1}+\dfrac1{\bar z_2}+\dfrac1{\bar z_3}\right)\left(\dfrac4{z_1}+\dfrac1{z_2}+\dfrac1{z_3}\right)=9$}
\opt{D}{orthocenter and circumcentre of $\Delta ABP$ will coincide.}
\end{qoptions}
\sect{SEC--2}
\paraq{1}{2}
A and B are playing a badminton match with the agreement that winner of each set will get 1 point and the loser 0 point. The match ends as soon as one of them is ahead by 2 points or number of sets reaches six. It is supposed that the probabilities of A and B winning a set are $\dfrac23$ and $\dfrac13$ respectively and each set is independent.\\
Let $X_i$ denotes the event that atleast $i$ sets are played and $Y$ and $Z$ denotes the event that match has won by A and B respectively.\\
On the basis of above information, answer the following questions :
\qn{1} Identify incorrect option --
\begin{qoptions}
\opt{A}{$P\left(\dfrac Y{X_5}\right)=\dfrac49$}
\opt{B}{$P\left(\dfrac Z{X_4}\right)=\dfrac{13}{81}$}
\opt{C}{$P(X_{2k-1})=P(X_{2k})\ \forall k\in\{1,2,3\}$}
\opt{D}{$P\left(\dfrac Z{X_1}\right)=\dfrac{64}{729}$}
\end{qoptions}
\qn{2} If it is known that A has won the third set then probability that A will win the match is equal to
\begin{qoptions}
\opt{A}{$\dfrac49$}
\opt{B}{$\dfrac{13}{18}$}
\opt{C}{$\dfrac{22}{27}$}
\opt{D}{$\dfrac{76}{81}$}
\end{qoptions}
\paraq{3}{4}
Consider the area $S_n$ bounded by $x$-axis, $y=e^{-x}\sin x$ between $x=n\pi$ and $x=(n+1)\pi$, where $x\ge0$\\
On the basis of above information, answer the following questions :
\qn{3} The value of $S_0$, is-
\begin{qoptions}
\opt{A}{$\dfrac12\left(1+e^{\pi}\right)$ sq. units}
\opt{B}{$\dfrac12\left(1+e^{-\pi}\right)$ sq. units}
\opt{C}{$\dfrac12\left(1-e^{-\pi}\right)$ sq. units}
\opt{D}{$\dfrac12\left(e^{\pi}-1\right)$ sq. units}
\end{qoptions}
\qn{4} $\displaystyle\sum_{n=0}^{\infty}S_n$ is equal to-
\begin{qoptions}
\opt{A}{$\dfrac{1+e^{\pi}}{1-e^{-\pi}}$}
\opt{B}{$\dfrac12\dfrac{\left(1+e^{-\pi}\right)}{\left(1-e^{-\pi}\right)}$}
\opt{C}{$\dfrac1{2\left(1-e^{-\pi}\right)}$}
\opt{D}{none of these}
\end{qoptions}
\paraq{5}{6}
Let $M_n=[M_{ij}]$ be a $n\times n$ matrix such that
\[M_{ij}=\begin{cases}a,&i=j\\1,&|i-j|=1\\0,&\text{otherwise}\end{cases}\]
Let $D_n$ denotes the determinant value of matrix $M_n$.\\
On the basis of above information, answer the following questions :
\qn{5} If $a=2$ then unit's place digit of $D_{2017}$ is
\begin{qoptions}
\opt{A}{5}
\opt{B}{6}
\opt{C}{7}
\opt{D}{8}
\end{qoptions}
\qn{6} If $a=1$ then $\displaystyle\sum_{k=1}^{2017}|D_k|$ is ? (where $|\ |$ denotes absolute value function)
\begin{qoptions}
\opt{A}{672}
\opt{B}{1007}
\opt{C}{1008}
\opt{D}{1345}
\end{qoptions}
\sect{SEC--3}
\qn{1} If $(1+ax+bx^2)^4=a_0+a_1x+a_2x^2+\ldots+a_8x^8$ where $a,b,a_0,a_1,\ldots,a_8\in R$ and $a_0+a_1+a_2\ne0$ and $\begin{vmatrix}a_0&a_1&a_2\\a_1&a_2&a_0\\a_2&a_0&a_1\end{vmatrix}=0$, then value of $5\left(\dfrac ab\right)$.
\qn{2} $z_1=(8+i)\sin\theta+(7+4i)\cos\theta$ and $z_2=(1+8i)\sin\theta+(4+7i)\cos\theta$. If $z_1.z_2=x+iy$ where $x,y\in R$ and $M$ is the greatest value of $x+y$ then value of $M^{1/3}$ ?
\qn{3} If the locus of middle point of contact of tangent drawn to the parabola $y^2=8x$ and the foot of perpendicular drawn from its focus to the tangent is a conic then find latus rectum of the conic.
\qn{4} Let $A\equiv(-4,0)$ and $B\equiv(4,0)$ if the number of points $C$ on the circle $x^2+y^2=16$, such that the area of triangle whose vertices are $A$, $B$ and $C$ is positive integer is $N$ then value of $\left[\dfrac N8\right]$ is equal to ?\\
(where [ ] denotes the greatest integer function)
\qn{5} A differentiable function $f$ satisfies the relation $f(x+y)=f(x)+f(y)+2xy(x+y)-\dfrac13\ \forall\ x,y\in R$ and $\lim\limits_{h\to0}\left(\dfrac{3f(h)-1}{6h}\right)=\dfrac23$, then value of $[f(2)]$ is, where [.] denote greatest integer function
\qn{6} Let $f(x)$ be a non constant thrice differentiable function defined on $(-\infty,\infty)$ such that $f(x)=f(6-x)$ and $f'(0)=0=f'(2)=f'(5)$. If $n$ is the minimum number of roots of $(f''(x))^2+f'(x)f'''(x)=0$ in the interval $[0,6]$ then value of $\dfrac n3$ is equal to ?
\qn{7} A curve is represented by $x=\sec^2t$ and $y=\cot t$, where $t$ is the parameter. If the tangent at point $P$ on curve where $t=\dfrac\pi4$ meets the curve again at $Q$ then $|PQ|=\dfrac{p\sqrt q}2$, where $p,q$ are coprime. Find $p+q$
\qn{8} $f(x)+f(y)=f\left(\dfrac{x+y}{1-xy}\right)$ for all $x,y\in R$, $(xy\ne1)$ and $\lim\limits_{x\to0}\dfrac{f(x)}x=2$. Then $f'(1)$ is
\qn{9} If the eccentricity of curve for which tangent at point $P$ intersects the $y$ axis at $M$ above $x$-axis such that origin is equidistant from $M$ and point of tangency is $e$ then value of $6e^2$ is?
\qn{10} The solution of $x^2\dfrac{dy}{dx}-xy=1+\cos\dfrac yx$ is $\tan\left(\dfrac y{px}\right)=c-\dfrac1{qx^2}$ where $c$ is constant then, find the value of $p+q$ ?

\end{multicols}
\end{document}
